\section{Related Work}
\label{sec:related-work}

Motion-sickness research commonly relates discomfort to conflicts among visual, vestibular, and proprioceptive motion information. Peripheral and anticipatory visual displays have consequently been studied as a way to convey vehicle dynamics while a passenger performs a non-driving task. Karjanto et al. evaluated a peripheral visual feedforward system in automated-driving scenarios, while Li et al. described a mobile motion-cueing application for non-driving tasks~\cite{karjanto2018feedforward,li2023mosi}.

The present work differs in its system boundary. The visual intervention is rendered on a Windows laptop rather than only on a phone or dedicated in-vehicle display, while the motion signal comes from a separate Android or iPhone companion. This separation makes transport latency, platform normalization, mount calibration, and stale-data behavior part of the intervention rather than incidental implementation details.

The user study uses the Fast Motion Sickness scale for repeated symptom reports and established questionnaires for susceptibility and multidimensional symptoms~\cite{keshavarz2011fms,golding2006mssq,gianaros2001msaq}. The final related-work review will be expanded after all candidate papers have been verified against publisher metadata.
